\documentclass[a4paper]{article}

\usepackage{graphicx}
\usepackage{amsmath,amssymb}
\usepackage{minted}
\usepackage{multirow}
\usepackage{float}
\usepackage{todonotes}
\usepackage{forest}
\usepackage{xstring}
\usepackage{xcolor}
\usepackage[most]{tcolorbox}
\usepackage{xparse}
\usepackage{natbib}
\usepackage{tikz}

\usetikzlibrary{graphs,graphdrawing}
\usegdlibrary{layered}

% for external graphviz
\input{/home/donkarlo/Dropbox/repo/nd_ai_project/src/nd_ai/knoweldge_representation/language/formal/computer/programming/tex/latex/code_clips/external_graphviz.tex}

% for tags
\input{/home/donkarlo/Dropbox/repo/nd_ai_project/src/nd_ai/knoweldge_representation/language/formal/computer/programming/tex/latex/parts/control_sequence/kind/semtag/semtag.tex}

% for links
\input{/home/donkarlo/Dropbox/repo/nd_ai_project/src/nd_ai/knoweldge_representation/language/formal/computer/programming/tex/latex/code_clips/seehere.tex}

\title{Validating autoencoder performance}
\date{}
\author{Mohammad Rahmani}

\begin{document}
    \maketitle
    
    
    \section{Validation of Autoencoder Performance}
        
        An autoencoder is usually validated by measuring how well it reconstructs unseen data.
        Let the validation dataset be
        
        \[
            X_{\mathrm{val}}=\{x_1,x_2,\dots,x_N\},
            \qquad
            x_i \in \mathbb{R}^{d}.
        \]
        
        For each input sample \(x_i\), the encoder maps the input into a latent representation,
        
        \[
            z_i=f_{\phi}(x_i),
        \]
        
        and the decoder reconstructs the input as
        
        \[
            \hat{x}_i=g_{\theta}(z_i)=g_{\theta}(f_{\phi}(x_i)).
        \]
        
        The main validation criterion is the reconstruction error on unseen validation data:
        
        \[
            \mathcal{L}_{\mathrm{val}}
            =
            \frac{1}{Nd}
            \sum_{i=1}^{N}
            \sum_{j=1}^{d}
            (x_{ij}-\hat{x}_{ij})^2.
        \]
        
        A lower validation reconstruction error indicates that the autoencoder reconstructs unseen samples more accurately.
        However, the training error must also be compared with the validation error:
        
        \[
            \Delta_{\mathrm{generalization}}
            =
            \mathcal{L}_{\mathrm{val}}
            -
            \mathcal{L}_{\mathrm{train}}.
        \]
        
        If this gap is small, the autoencoder is likely generalizing well.
        If the training error is low but the validation error is high, the model is probably overfitting.
    
    
    \section{Common Reconstruction Metrics}
        
        Several metrics can be used to evaluate the reconstruction quality.
        
        Mean squared error is defined as
        
        \[
            \mathrm{MSE}
            =
            \frac{1}{Nd}
            \sum_{i=1}^{N}
            \sum_{j=1}^{d}
            (x_{ij}-\hat{x}_{ij})^2.
        \]
        
        Mean absolute error is defined as
        
        \[
            \mathrm{MAE}
            =
            \frac{1}{Nd}
            \sum_{i=1}^{N}
            \sum_{j=1}^{d}
            \left|x_{ij}-\hat{x}_{ij}\right|.
        \]
        
        For multidimensional sensory data, the reconstruction error of each sample can be computed as
        
        \[
            e_i
            =
            \frac{1}{d}
            \sum_{j=1}^{d}
            (x_{ij}-\hat{x}_{ij})^2,
        \]
        
        where \(d\) is the number of input dimensions.
    
    
    \section{Validation for Novelty or Anomaly Detection}
        
        When an autoencoder is used for novelty detection, it is usually trained only on normal data.
        The assumption is that normal samples are reconstructed well, while novel or abnormal samples produce larger reconstruction errors.
        
        The reconstruction error for a validation sample is
        
        \[
            e_i
            =
            \frac{1}{d}
            \sum_{j=1}^{d}
            (x_{ij}-\hat{x}_{ij})^2.
        \]
        
        A threshold \(\tau\) can be selected from the reconstruction errors of normal validation data.
        For example, using the \(q\)-quantile:
        
        \[
            \tau
            =
            Q_q
            (
            \{e_1,e_2,\dots,e_N\}
            ).
        \]
        
        A new sample \(x_i\) is then classified as novel if
        
        \[
            e_i > \tau.
        \]
        
        Otherwise, it is classified as normal:
        
        \[
            e_i \leq \tau.
        \]
        
        If labeled normal and abnormal validation data are available, the threshold can be evaluated using precision, recall, \(F_1\)-score, ROC-AUC, or PR-AUC.
    
    
    \section{Validation of the Latent Space}
        
        Reconstruction error alone is not always sufficient.
        The latent space should also be inspected.
        A useful autoencoder should map similar inputs to nearby latent representations.
        
        For two input samples \(x_i\) and \(x_j\), their latent vectors are
        
        \[
            z_i=f_{\phi}(x_i),
            \qquad
            z_j=f_{\phi}(x_j).
        \]
        
        The latent distance can be computed as
        
        \[
            d_z(i,j)
            =
            \|z_i-z_j\|_2.
        \]
        
        If similar input samples have small latent distances and different input samples have larger latent distances, the latent representation is more meaningful.
    
    
    \section{Final Test Evaluation}
        
        The validation set is used for model selection and threshold selection.
        The test set must be kept independent and used only once for the final evaluation.
        
        The final test loss is
        
        \[
            \mathcal{L}_{\mathrm{test}}
            =
            \frac{1}{N_{\mathrm{test}}d}
            \sum_{i=1}^{N_{\mathrm{test}}}
            \sum_{j=1}^{d}
            (x_{ij}^{\mathrm{test}}-\hat{x}_{ij}^{\mathrm{test}})^2.
        \]
        
        The autoencoder is considered reliable if the validation loss and test loss are close:
        
        \[
            \mathcal{L}_{\mathrm{test}}
            \approx
            \mathcal{L}_{\mathrm{val}}.
        \]
        
        A large difference between test loss and validation loss indicates that the validation procedure may not represent the true generalization behavior of the model.
    
    
    \section{Practical Validation Procedure}
        
        The practical validation of an autoencoder can be summarized as follows:
        
        \begin{enumerate}
            \item Split the data into training, validation, and test sets.
            \item Train the autoencoder only on the training set.
            \item Compute reconstruction error on the validation set.
            \item Compare training and validation losses to check overfitting.
            \item Inspect reconstructed samples visually or numerically.
            \item Analyze the latent space using clustering or dimensionality reduction.
            \item If the autoencoder is used for novelty detection, choose a threshold on validation errors.
            \item Evaluate the final performance on the independent test set.
        \end{enumerate}
    
    
    \section{Interpretation}
        
        An autoencoder is considered to perform well if it satisfies the following conditions:
        
        \begin{enumerate}
            \item The validation reconstruction error is low.
            \item The gap between training and validation error is small.
            \item The reconstructed signals preserve the important structure of the original signals.
            \item The latent space separates meaningful patterns.
            \item For novelty detection, abnormal samples produce higher reconstruction errors than normal samples.
            \item The final test loss is close to the validation loss.
        \end{enumerate}

%	\bibliographystyle{plainnat}
%	\bibliography{/home/donkarlo/Dropbox/repo/nd_shared_research/bibliography/bibliography.bib}
\end{document}